\documentclass{article}
\usepackage{amsmath}  % For advanced math
\usepackage{amssymb}  % For \mathbb

\begin{document}

\section*{Complex Numbers}

\begin{itemize}
    \item \textbf{Imaginary Unit}: The imaginary unit is defined as:
        \[
        i = \sqrt{-1}, \quad \text{so that } i^2 = -1
        \]
    
    \item \textbf{Complex Number}: A complex number is any number of the form:
        \[
        z = a + bi
        \]
        where:
        \begin{itemize}
            \item \( a \) is the real part (a real number).
            \item \( bi \) is the imaginary part (a real number multiplied by \( i \)).
        \end{itemize}
    
    \item \textbf{Pure Imaginary Number}: If \( a = 0 \), then the complex number \( z = bi \) is called a \textit{pure imaginary number}.
    
    \item \textbf{Addition of Complex Numbers}:
        \[
        (a + bi) + (c + di) = (a + c) + (b + d)i
        \]
    
    \item \textbf{Multiplication of Complex Numbers}:
        \[
        (a + bi) \cdot (c + di) = (ac - bd) + (ad + bc)i
        \]
    
    \item \textbf{Conjugate of a Complex Number}: The conjugate of a complex number \( z = a + bi \) is denoted as \( \overline{z} \) and is defined as:
        \[
        \overline{z} = a - bi
        \]
    
    \item \textbf{Modulus (Magnitude) of a Complex Number}: The modulus or magnitude of a complex number \( z = a + bi \) is given by:
        \[
        |z| = \sqrt{a^2 + b^2}
        \]
    
    \item \textbf{Division of Complex Numbers}:
        \[
        \frac{a + bi}{c + di} = \frac{(a + bi) \cdot (c - di)}{(c + di) \cdot (c - di)} = \frac{(a + bi) \cdot (c - di)}{c^2 + d^2}
        \]
\end{itemize}

\end{document}
